\subsection{AI in Urban Planning}

Cities are dynamic systems influenced by a wide range of interacting factors.
In large urban environments, multifunctional activities often coexist, and their interactions causes complex and conflicting urban planning challenges.

Large language models can serve as intelligent assistants for urban planners, supporting the synthesis of conceptual ideas, generation of urban design proposals, and evaluation of planning outcomes \cite{zheng2025urban}.
Multi-agent LLM approaches further enable effective handling of complex urban planning tasks while maintaining high relevance of generated outputs, making them practical for real-world applications in urban governance \cite{kalyuzhnaya2025llm}.

Existing digital planning solutions often struggle to adapt to the rapidly changing nature of urban environments and are typically designed for single datasets or specific use cases, which limits their generalizability across different planning tasks \cite{XIA2025106142}.
In \cite{singla2024adaptive}, agent-based approaches are explored for urban planning, focusing on optimizing basic infrastructure requirements at the neighborhood level to achieve balanced and coherent urban development.
The authors of \cite{jiao2024city} aim toward more transparent city management.
Their approach leverages pretrained end-to-end optimization models and stakeholder-oriented natural language query features to ensure robust performance under uncertainty. PlanGPT \cite{zhu2024plangpt} successfully addresses challenges in urban planning document management, review, and assessment, significantly improving timeliness and multimodal handling of planning artifacts.
Incorporating demographic and life-value data, the work in \cite{gao2025simulating} proposes a predictive framework for anticipating stakeholder responses via simulation-based modeling, thereby enabling more equitable and cost-effective urban planning decisions.
The study in \cite{LIU2025100666} investigates transparent AI deployment in urban planning and establishes systematic design principles supporting sustainable urban development through effective human–AI collaboration.

Traditional approaches rely heavily on extensive historical data and assume relatively stable relationships between variables, whereas LLM-based systems demonstrate strong performance even in low-shot or zero-shot forecasting scenarios by leveraging deeper contextual understanding and interdisciplinary knowledge \cite{pan2025tools}.
However, LLMs often struggle with heterogeneous data streams originating from multiple sources, especially when these span different spatial resolutions and temporal scales \cite{luo2026geojson}.
To address this challenge, CityGPT \cite{feng2025citygpt} introduces a systematic framework for evaluating and enhancing the spatial cognition abilities of general-purpose LLMs in urban environments.
Another approach, ShapefileGPT \cite{lin2025shapefilegpt}, effectively processes complex vector geospatial data, demonstrating strong spatial reasoning and analytical capabilities.

In \cite{sung5382891planning}, a comparison is conducted between human-generated plans and ComPlanAI, a hybrid AI system combining rule-based computation with large language models for comprehensive urban planning.
The results highlight the potential of AI-assisted co-design in strategic planning contexts.
However, more comprehensive evaluation of planning quality requires integration of cartographic and GIS-based data into assessment frameworks.
The study in \cite{luo2026geojson} addresses this limitation by transforming natural language instructions into structured GeoJSON operations using LLM enhancement techniques such as function calling and code generation.